\subsection{Positive quadratic-drift estimates}
\label{sec:direct-positive-quadratic}

Each branch on this chart has a first boundary root. When the central
coefficient keeps these roots away from the central line, we compare
the endpoint slopes either by curvature or, near the symmetric limit,
by distances to the roots. A further rescaling treats the root that
approaches the central line.

Suppose
\(q_2=s>0,\ q_0<0,\ q_1>0\), and write
\begin{equation}
\label{eq:direct-positive-secondary-chart}
 r=\max\{\sqrt{|q_0|/s^3},q_1/s^2\},\quad
 q_0=-as^3r^2,\quad q_1=bs^2r,\quad
 \max\{\sqrt a,b\}=1 .
\end{equation}
The small parameters \(s,r>0\) are independent. Put
\(t_0=s^2r^2,\ Q=st_0\), and \(X=|\eta|/(sr)\).
Here \(sr\) is the radius scale, \(Q=s^3r^2\) the central
height scale, and \(t_0\) the fixed-output height scale;
the common logarithmic height coordinate used below is \(T=-s\log h\).
The outward drifts are
\begin{equation}
\label{eq:direct-positive-secondary-drifts}
 K_\sigma(X)=a-\sigma bX-X^2
                 -\sigma rX^3R_q(\sigma srX),\qquad
 F_\sigma=K_\sigma/X .
\end{equation}
\begin{proposition}[Positive quadratic drift]
\label{prop:direct-positive-quadratic}
There is \(\kappa_+>0\) such that, for all sufficiently small
independent \(s,r>0\) and every \(L>1/\kappa_+\), the return
displacement has at most three distinct zeros on the part of
\eqref{eq:direct-positive-secondary-chart} outside the strip
\(|q_0|\le Lq_1^2\).
\end{proposition}

\subsubsection{Constant coefficient bounded away from zero}

First fix \(a\ge a_0^+>0\). This lower bound keeps both roots in a
compact positive radius interval, where we can prepare the endpoint
maps uniformly. The roots are simple and satisfy
\[
 \alpha_\sigma
 =\frac{\sqrt{b^2+4a}-\sigma b}{2}+O(r),\qquad
 c\le\alpha_\sigma\le C,\qquad
 K_\sigma'(\alpha_\sigma)<-c .
\]
At height \(z=t_0Y\),
\[
 X_{\log Y}=sF_\sigma(X)+YB(\sigma srX,t_0Y)/X .
\]
The boundary-flow coordinate applies with \(\delta=s\),
since the error in \(B\) is \(O(sr)=O(s)\).
At a fixed small output height \(Y_f>0\) the coefficient ratios
are \(O(s)\), and \cref{lem:boundary-flow-endpoint} with
\(\ell=sr\) and \(\delta=s\) gives
\begin{equation}
\label{eq:direct-positive-secondary-endpoints}
 -I_{\sigma,w}=EA_\sigma wV_\sigma(w),\qquad
 E=t_0(z_*/(t_0Y_f))^{-2s}\to0,\qquad
 V_\sigma=1+O_{C^3}(s),
\end{equation}
on fixed intervals of positive projected states up to the roots.
Here \(A_\sigma\asymp1\) and
\(A_-/A_+\le\exp(2C_{\rm cone}\sqrt{z_*}+Cs)\).

Use the shifted coordinate \(T_Q=T+s\log Q\), as in \eqref{eq:direct-source-clock-convention}.
The phases are
\begin{equation}
\label{eq:direct-positive-secondary-phases}
 \begin{aligned}
 w_\sigma&=\Phi_{F_\sigma}^{\,T_Q+C_\sigma+
                                  s\Xi_\sigma(\xi_\sigma)}(x_{L,\sigma}),
 &T_Q&=-s\log(h/Q),\\
 \xi_\sigma&=\e^{\mathcal A_\sigma(x_{L,\sigma})/s}h/Q,
 &C_\sigma&=s\log(Y_f/s)-\mathcal A_\sigma(x_{L,\sigma}),\\
 \mathcal A_\sigma(x)&=\int_0^x u/K_\sigma(u)\,du .
 \end{aligned}
\end{equation}
Choose both regular entry actions below one quarter of the minimum
positive terminal action. Then \(w_+\ge\alpha_+/2\) forces a
fixed advance on both legs. The negative state lies in a fixed
interval away from zero; by \cref{rem:direct-projected-continuation}
it stays below \(\alpha_-\), and the inverse projection continues the
negative graph.

Suppose also \(b\ge \underline b>0\). The drifts
\eqref{eq:direct-positive-secondary-drifts} have
\(K_\sigma'(0)=-\sigma b\) and, in the notation of
\eqref{eq:direct-curvature-coefficient-cubic},
\(R_{K_\sigma}(w)=-\sigma rR_q(\sigma srw)=O_{C^2}(r)\). Hence
\eqref{eq:direct-curvature-coefficient-cubic} gives, on these compact
intervals,
\[
 \frac{wK_+''-K_+'}{w^2}=\frac b{w^2}+O(r)>c,\qquad
 \frac{wK_-''-K_-'}{w^2}=-\frac b{w^2}+O(r)<-c .
\]
Since the negative state stays in a fixed interval away from zero,
\cref{lem:curvature-criterion} gives \(\mathcal L''<0\) on this root
interval, including all approaches to its ends.

To extend this comparison to larger central heights, join the root
interval to the regular prefix.
Apply \cref{lem:prefix-root-cover} with \(\delta=s\), \(\ell=sr\),
\(\epsilon_3=r\), comparison drifts \(a-\sigma bx-x^2\),
\(R=3\alpha_+/4\) and \(w_a=\alpha_+/2\).
The positive drift margin and the gap \(R-w_a\) are fixed, and the
boundary-flow inverse is defined up to \(\alpha_+\) at \(Y_f\).
Since \(\mathcal L_{TT}<0\) on this root range, the lemma gives at most
three zeros of \(\Delta_p\) on \(\cH_p\).

\subsubsection{Small linear coefficient: comparison in root-distance coordinates}

For small \(b\), the curvature comparison loses its fixed margin.
We use linearizing coordinates for the two scalar boundary flows to
obtain a power relation between the distances to the roots, then
combine it with the exterior contraction. Put
\[
 k_\sigma=-F_\sigma'(\alpha_\sigma)
                  =2+O(b+r)>c .
\]
There is a positive root-distance coordinate
\(\ell_\sigma(w)\), decreasing to zero, such that
\begin{equation}
\label{eq:direct-positive-root-distance}
 F_\sigma(w)\ell_\sigma'(w)=-k_\sigma\ell_\sigma(w),
 \qquad \ell_\sigma'(\alpha_\sigma)=-\alpha_\sigma .
\end{equation}
To construct this coordinate, factor out \(\alpha_\sigma-w\).
The logarithmic derivative of the remaining unit is
\(-k_\sigma/F_\sigma+1/(\alpha_\sigma-w)\), whose pole
is removable. Integration and normalization at the root
give uniform derivative bounds and a regular inverse.
Extend to the regular entry using the same scalar equation.
At \(b=r=0\) the formulas are
\[
 F_\sigma=(a-w^2)/w,\qquad k_\sigma=2,\qquad
 \ell_\sigma=(a-w^2)/2 .
\]

The constants
\[
 c_\sigma^{\rm lin}=\ell_\sigma(x_{L,\sigma})
              \exp[k_\sigma \mathcal A_\sigma(x_{L,\sigma})]
\]
are independent of the entry chosen, by
\eqref{eq:direct-positive-root-distance}.
They have positive bounded values and smooth dependence on
the compact coefficients. The phase identities give
\begin{equation}
\label{eq:direct-positive-root-distance-phase}
 \ell_\sigma(w_\sigma)
 =c_\sigma^{\rm lin}(Y_f/s)^{-sk_\sigma}
       \e^{-k_\sigma T_Q}\e^{-sk_\sigma\Xi_\sigma(\xi_\sigma)} .
\end{equation}
Put \(x=\ell_+(w_+)\), \(y=\ell_-(w_-)\), and
\(\nu_k=k_-/k_+\). Since \(\nu_k k_+=k_-\),
\[
 \frac{y}{x^{\nu_k}}
 =\frac{c_-^{\rm lin}}{(c_+^{\rm lin})^{\nu_k}}
   \exp\!\left(-sk_-[\Xi_-(\xi_-)-\Xi_+(\xi_+)]\right).
\]
The exponential factors involving \(T_Q\) and \(Y_f/s\) cancel, giving
\begin{equation}
\label{eq:direct-positive-two-root-power}
 y=u_px^{\nu_k}(1+\rho_p(x)),\qquad
 \nu_k=k_-/k_+=1+O(b+r),\qquad
 u_p=c_-^{\rm lin}/(c_+^{\rm lin})^{\nu_k}=1+O(b+r).
\end{equation}
At \(b=r=0\), \(c_+^{\rm lin}=c_-^{\rm lin}=a/2\) and
\(\nu_k=1\), so the leading unit is exactly one. Its estimate
is uniform because \(a\ge a_0^+\).

On a sufficiently small fixed collar \(0<x<x_c\), a fixed root
distance forces a positive phase gap on both legs. The argument of
\cref{lem:direct-weighted-phase}, now using
\(F_\sigma\asymp\ell_\sigma\) near the simple roots, gives
\[
 |\partial_T^j\xi_\sigma|
 \le\varepsilon_{N,j}(s)\ell_\sigma(w_\sigma)^N,
 \qquad\varepsilon_{N,j}(s)\to0
\]
for every fixed \(N,j\). The two distances are comparable through fixed
positive powers by \eqref{eq:direct-positive-root-distance-phase}, and
\(d\log x/dT=-k_+\partial_T\tau_+\) stays bounded away from zero.
Thus substitution in \eqref{eq:direct-positive-two-root-power} gives
\[
 |(x\partial_x)^j\rho_p(x)|\le C_jx^\vartheta,
 \qquad 0\le j\le2,
\]
with fixed \(\vartheta>0\), uniformly as \(h\downarrow0\).

In these root-distance coordinates the endpoints are
\[
 I_\sigma(\ell_\sigma)=b_\sigma+E\mathcal F_\sigma(\ell_\sigma),\qquad
 \mathcal F_\sigma(0)=0,\quad
 c\le\mathcal F_\sigma'\le C,\quad
 \|\mathcal F_\sigma\|_{C^3}\le C .
\]
The normalization
\(\ell_\sigma'(\alpha_\sigma)=-\alpha_\sigma\) and
\eqref{eq:direct-positive-secondary-endpoints} give
\[
 \mathcal F_-'(0)/\mathcal F_+'(0)
                   \le\exp(2C_{\rm cone}\sqrt{z_*}+C(s+b+r)).
\]
If some central height in the collar lies in \(\cH_p\), the endpoints
on the collar, including the reference endpoints \(b_\sigma\), differ
from admissible endpoints by \(O(E)\), so they lie in \(\widetilde W_\pm\)
by \cref{lem:exterior-multiplier-new}.
In these coordinates, the logarithmic slope ratio is
\[
 \mathcal L(x)
 =\log\left[\Pi_p'(I_-(y(x)))
       \frac{\mathcal F_-'(y(x))y'(x)}{\mathcal F_+'(x)}\right]
 =c_p+(\nu_k-1)\log x+e_p(x),
\]
where
\[
 c_p=\log\left[\Pi_p'(b_-)
        \frac{\mathcal F_-'(0)}{\mathcal F_+'(0)}u_p\nu_k\right]
       <-c_0,\qquad |e_p|+|xe_p'|\le Cx^\vartheta.
\]
Choose \(b_{*,+}\), then \(r,s\), for this strict margin.
The nonlinear exterior contributes \(O(Ey(x))\) with
the same Euler estimate. Since
\(dx/dh=I_{+,h}/(E\mathcal F_+'(x))>0\), the common factor
\(dx/dh\) cancels from the displayed quotient, and \(x_T<0\).
\Cref{lem:one-logarithm-zero-count} shows that every zero in a
sufficiently small fixed collar satisfies \(x\mathcal L_x<0\), and that
such zeros occur only if \(\nu_k<1\). Consequently
\(\mathcal L_T>0\) at every collar zero.

Choose \(R=(1-\epsilon_R)\alpha_+\), with fixed small \(\epsilon_R>0\), strictly inside
the prepared collar. The positive drift margin is fixed. Apply
\cref{lem:prefix-root-cover} on the same scales, choosing \(w_a<R\)
inside this collar with a fixed gap. Every collar zero has \(\mathcal L_T>0\), so
\(\Delta_p\) has at most two zeros for \(0\le b\le b_{*,+}\).
Taking \(b\ge b_{*,+}/2\) in the curvature
argument gives an overlap and covers the compact \(a\) region.

\subsubsection{The central weak root}
\label{sec:positive-quadratic-weak-root}

It remains to allow the constant coefficient to approach zero. For
small \(a>0\), normalization in
\eqref{eq:direct-positive-secondary-chart} forces \(b=1\).
Put \(\kappa=s/a\). With \(\ell=sra\), the scales
\eqref{eq:weak-root-scales} hold with \(\hat r=r\) and \(\epsilon_2=-a\).
At the radius \(\ell_d=srd\) the negative drift is
\[
 K_d(x)=a_d+x-dx^2+rd^2x^3R_q(-\ell_dx),
\]
so \eqref{eq:weak-root-negative-drift} holds for small fixed
\(d_{\max}=d_0\) and \(r\le r(d_0)\).

For \eqref{eq:weak-root-outer}, consider the negative graphs with
\(z<Y_b|\eta|^2\) at \(|\eta|=srd_0\). In the secondary radius
\(X=|\eta|/(sr)\) the negative drift is
\[
 K_-(X)=a+X-X^2+rX^3R_q(-srX),
\]
which is the drift \(K_1\) of \cref{cor:weak-root-negative-first-root},
since \(\ell/a=sr\). Its first root is \(\beta=1+O(a+r)\le2\), and
\(K_-'/X\ge-C\) on \([d_0/2,\beta]\).
\Cref{cor:weak-root-negative-first-root}, with \(d_{\max}=d_0\) and
\(\beta^*=2\), gives \(\mathfrak H_-\ge-C-1\) and \(|I_{-,T}|=o(1)\) for
these graphs, up to the limiting root. \Cref{prop:weak-root-passage} gives
at most two zeros of \(\Delta_p\) for \(0<a\le a_+\) and
\(0<\kappa\le\kappa_+\), where \(a_+\) and \(\kappa_+\) are its constants
\(a_\circ\) and \(\kappa_\circ\).

The complement of \(\kappa\le\kappa_+\) in the region \(0<a\le a_+\) has
\(|q_0|/q_1^2=a/s<1/\kappa_+\), and lies in the strip
\(|q_0|\le Lq_1^2\) for every \(L>1/\kappa_+\). The compact-\(a\)
argument, applied down to a threshold \(a_0^+<a_+\), overlaps the
weak-root region. The curvature, root-distance, and weak-root
comparisons give bounds three, two, and two, respectively, so their
maximum is three outside that strip. The thresholds for \(s,r\) are those
of the three regions and are independent of \(L\).
This proves \cref{prop:direct-positive-quadratic}.
